% Chants of the Missa pro defunctis, the absolution and In paradisum:
% \input by missa-pro-defunctis.tex.  Score paths are relative to a sibling
% directory, so the file also works when \input from another booklet.
% ---------------------------------------------------------------------------
\chant{Requiem aeternam}{introitus · tonus VI}
\lilypondfile{../missa-pro-defunctis/scores/introitus.ly}
\begin{tekst}
\w{Réquiem aetérnam dona eis, Dómine: et lux perpétua lúceat eis.}
  {Wieczny odpoczynek racz im dać, Panie, a światłość wiekuista niechaj im świeci.}
\w{\rubr{Ps. 64.} Te decet hymnus, Deus, in Sion, et tibi reddétur votum in Jerúsalem:
   exáudi oratiónem meam, ad te omnis caro véniet.}
  {\rubr{Ps. 65.} Tobie, Boże, przystoi hymn na Syjonie, Tobie składać się będą śluby
   w Jeruzalem; wysłuchaj mojej modlitwy, do Ciebie przyjdzie wszelkie ciało.}
\w{\rubr{Powtarza się antyfonę} Réquiem aetérnam.}
  {\rubr{Powtarza się antyfonę} Wieczny odpoczynek.}
\end{tekst}

% ---------------------------------------------------------------------------
\chant{Kyrie}{tonus VI}
\lilypondfile{../missa-pro-defunctis/scores/kyrie.ly}
\begin{tekst}
\w{Kýrie, eléison. Christe, eléison. Kýrie, eléison.}
  {Panie, zmiłuj się nad nami. Chryste, zmiłuj się nad nami.
   Panie, zmiłuj się nad nami.}
\end{tekst}

\clearpage
% ---------------------------------------------------------------------------
\chant{Requiem aeternam}{graduale · tonus II}
\lilypondfile{../missa-pro-defunctis/scores/graduale.ly}
\begin{tekst}
\w{Réquiem aetérnam dona eis, Dómine: et lux perpétua lúceat eis.}
  {Wieczny odpoczynek racz im dać, Panie, a światłość wiekuista niechaj im świeci.}
\w{\V In memória aetérna erit justus: ab auditióne mala non timébit.}
  {\V Sprawiedliwy będzie w wiecznej pamięci, nie ulęknie się złej wieści.}
\end{tekst}

% ---------------------------------------------------------------------------
\chant{Absolve, Domine}{tractus · tonus VIII}
\lilypondfile{../missa-pro-defunctis/scores/tractus.ly}
\begin{tekst}
\w{Absólve, Dómine, ánimas ómnium fidélium defunctórum ab omni vínculo delictórum.}
  {Uwolnij, Panie, dusze wszystkich wiernych zmarłych od wszelkich więzów grzechu.}
\w{\V Et grátia tua illis succurrénte, mereántur evádere judícium ultiónis.}
  {\V I niech za wsparciem Twojej łaski zasłużą na uniknięcie wyroku potępienia.}
\w{\V Et lucis aetérnae beatitúdine pérfrui.}
  {\V I niech cieszą się szczęściem wiecznej światłości.}
\end{tekst}

% ---------------------------------------------------------------------------
\chant{Dies irae}{sequentia · tonus I}
\lilypondfile{../missa-pro-defunctis/scores/sequentia.ly}
\begin{tekst}
\w{Dies irae, dies illa, solvet saeclum in favílla: teste David cum Sibýlla.}
  {Dzień gniewu, ów dzień obróci świat w popiół, jak świadczą Dawid i Sybilla.}
\w{Quantus tremor est futúrus, quando judex est ventúrus, cuncta stricte discussúrus!}
  {Jakież będzie drżenie, gdy przyjdzie Sędzia, by wszystko surowo rozsądzić!}
\w{Tuba mirum spargens sonum per sepúlcra regiónum, coget omnes ante thronum.}
  {Trąba, rozsiewając przedziwny dźwięk po grobach wszystkich krain,
   zgromadzi wszystkich przed tronem.}
\w{Mors stupébit et natúra, cum resúrget creatúra, judicánti responsúra.}
  {Zdumieje się śmierć i natura, gdy powstanie stworzenie,
   by dać odpowiedź Sędziemu.}
\w{Liber scriptus proferétur, in quo totum continétur, unde mundus judicétur.}
  {Przyniesiona będzie spisana księga, w której zawarte jest wszystko,
   według czego świat będzie sądzony.}
\w{Judex ergo cum sedébit, quidquid latet apparébit: nil inúltum remanébit.}
  {Gdy więc Sędzia zasiądzie, wyjdzie na jaw wszystko, co ukryte,
   nic nie zostanie bez kary.}
\w{Quid sum miser tunc dictúrus? Quem patrónum rogatúrus, cum vix justus sit secúrus?}
  {Cóż wtedy powiem ja, nędzny? Kogo poproszę o obronę,
   skoro i sprawiedliwy ledwie będzie bezpieczny?}
\w{Rex treméndae majestátis, qui salvándos salvas gratis, salva me, fons pietátis.}
  {Królu straszliwego majestatu, który darmo zbawiasz tych, co mają być zbawieni,
   zbaw mnie, źródło litości.}
\w{Recordáre, Jesu pie, quod sum causa tuae viae: ne me perdas illa die.}
  {Pamiętaj, dobry Jezu, że jestem przyczyną Twojej drogi:
   nie gub mnie w owym dniu.}
\w{Quaerens me, sedísti lassus: redemísti crucem passus: tantus labor non sit cassus.}
  {Szukając mnie, usiadłeś znużony; odkupiłeś mnie, cierpiąc krzyż:
   niech tak wielki trud nie pójdzie na marne.}
\w{Juste judex ultiónis, donum fac remissiónis ante diem ratiónis.}
  {Sprawiedliwy Sędzio, który karzesz, udziel daru przebaczenia
   przed dniem rozrachunku.}
\w{Ingemísco, tamquam reus: culpa rubet vultus meus: supplicánti parce, Deus.}
  {Wzdycham jak winowajca, od winy płonie moje oblicze:
   oszczędź błagającego, Boże.}
\w{Qui Maríam absolvísti, et latrónem exaudísti, mihi quoque spem dedísti.}
  {Ty, który rozgrzeszyłeś Marię i wysłuchałeś łotra,
   i mnie dałeś nadzieję.}
\w{Preces meae non sunt dignae: sed tu bonus fac benígne, ne perénni cremer igne.}
  {Moje prośby nie są godne, lecz Ty, dobry, spraw łaskawie,
   abym nie spłonął w wiecznym ogniu.}
\w{Inter oves locum praesta, et ab haedis me sequéstra, státuens in parte dextra.}
  {Daj mi miejsce wśród owiec i oddziel mnie od kozłów,
   stawiając po prawej stronie.}
\w{Confutátis maledíctis, flammis ácribus addíctis: voca me cum benedíctis.}
  {Gdy potępieni zostaną odrzuceni i wydani srogim płomieniom,
   wezwij mnie z błogosławionymi.}
\w{Oro supplex et acclínis, cor contrítum quasi cinis: gere curam mei finis.}
  {Proszę pokornie, pochylony, z sercem skruszonym jak popiół:
   miej w swej pieczy moją ostatnią godzinę.}
\w{Lacrimósa dies illa, qua resúrget ex favílla judicándus homo reus:
   huic ergo parce, Deus.}
  {Pełen łez będzie ów dzień, w którym powstanie z popiołu
   człowiek winny, by stanąć na sądzie: jemu więc przebacz, Boże.}
\w{Pie Jesu Dómine, dona eis réquiem. Amen.}
  {Dobry Jezu, Panie, daj im odpoczynek. Amen.}
\end{tekst}

% ---------------------------------------------------------------------------
\chant{Domine Jesu Christe}{offertorium · tonus II}
\lilypondfile{../missa-pro-defunctis/scores/offertorium.ly}
\begin{tekst}
\w{Dómine Jesu Christe, Rex glóriae, líbera ánimas ómnium fidélium
   defunctórum de poenis inférni, et de profúndo lacu:}
  {Panie Jezu Chryste, Królu chwały, wybaw dusze wszystkich wiernych
   zmarłych od kar piekła i od otchłani bez dna;}
\w{líbera eas de ore leónis, ne absórbeat eas tártarus, ne cadant in obscúrum:}
  {wybaw je z paszczy lwa, niech ich nie pochłonie piekło,
   niech nie wpadną w ciemności;}
\w{sed sígnifer sanctus Míchael repraeséntet eas in lucem sanctam:}
  {lecz chorąży, święty Michał, niech je stawi w światłości świętej,}
\w{Quam olim Abrahae promisísti, et sémini ejus.}
  {którą niegdyś obiecałeś Abrahamowi i jego potomstwu.}
\w{\V Hóstias et preces tibi, Dómine, laudis offérimus:
   tu súscipe pro animábus illis, quarum hódie memóriam fácimus:
   fac eas, Dómine, de morte transíre ad vitam.}
  {\V Ofiary i modły pochwalne Tobie, Panie, składamy:
   przyjmij je za te dusze, które dziś wspominamy;
   daj im, Panie, przejść od śmierci do życia,}
\w{\rubr{Powtarza się} Quam olim Abrahae \rubr{do końca.}}
  {\rubr{Powtarza się} którą niegdyś obiecałeś \rubr{do końca.}}
\end{tekst}

% ---------------------------------------------------------------------------
\chant{Sanctus}{jak w Mszy XVIII}
\lilypondfile{../missa-pro-defunctis/scores/sanctus.ly}
\begin{tekst}
\w{Sanctus, Sanctus, Sanctus Dóminus Deus Sábaoth.}
  {Święty, Święty, Święty, Pan Bóg Zastępów.}
\w{Pleni sunt caeli et terra glória tua.}
  {Pełne są niebiosa i ziemia chwały Twojej.}
\w{Hosánna in excélsis.}{Hosanna na wysokości.}
\w{Benedíctus qui venit in nómine Dómini.}
  {Błogosławiony, który idzie w imię Pańskie.}
\w{Hosánna in excélsis.}{Hosanna na wysokości.}
\end{tekst}

% ---------------------------------------------------------------------------
\chant{Agnus Dei}{jak w Mszy XVIII}
\lilypondfile{../missa-pro-defunctis/scores/agnus.ly}
\begin{tekst}
\w{Agnus Dei, qui tollis peccáta mundi: dona eis réquiem.}
  {Baranku Boży, który gładzisz grzechy świata, daj im odpoczynek.}
\w{Agnus Dei, qui tollis peccáta mundi: dona eis réquiem.}
  {Baranku Boży, który gładzisz grzechy świata, daj im odpoczynek.}
\w{Agnus Dei, qui tollis peccáta mundi: dona eis réquiem sempitérnam.}
  {Baranku Boży, który gładzisz grzechy świata, daj im wieczny odpoczynek.}
\end{tekst}

\clearpage
% ---------------------------------------------------------------------------
\chant{Lux aeterna}{communio · tonus VIII}
\lilypondfile{../missa-pro-defunctis/scores/communio.ly}
\begin{tekst}
\w{Lux aetérna lúceat eis, Dómine: Cum sanctis tuis in aetérnum, quia pius es.}
  {Światłość wiekuista niechaj im świeci, Panie, ze świętymi Twoimi na wieki,
   bo jesteś łaskawy.}
\w{\V Réquiem aetérnam dona eis, Dómine, et lux perpétua lúceat eis.}
  {\V Wieczny odpoczynek racz im dać, Panie, a światłość wiekuista niechaj im świeci.}
\w{Cum sanctis tuis in aetérnum, quia pius es.}
  {Ze świętymi Twoimi na wieki, bo jesteś łaskawy.}
\end{tekst}

% ---------------------------------------------------------------------------
\chant{Requiescant in pace}{zamiast Ite, missa est}
\lilypondfile{../missa-pro-defunctis/scores/requiescant.ly}
\begin{tekst}
\w{\V Requiéscant in pace.}{\V Niech odpoczywają w pokoju.}
\w{\R Amen.}{\R Amen.}
\end{tekst}

\clearpage
% ---------------------------------------------------------------------------
\begin{center}
  {\LARGE\scshape Absolutio}\par\smallskip
  {\large\itshape Obrzęd ostatniego pożegnania}\par\smallskip
  {\small\rubr{Po Mszy, przy trumnie lub katafalku.}}
\end{center}

\chant{Libera me, Domine}{responsorium · tonus I}
\lilypondfile{../missa-pro-defunctis/scores/libera.ly}
\begin{tekst}
\w{Líbera me, Dómine, de morte aetérna, in die illa treménda:
   Quando caeli movéndi sunt et terra:
   Dum véneris judicáre saeculum per ignem.}
  {Wybaw mnie, Panie, od śmierci wiecznej w ów dzień straszny,
   gdy poruszą się niebiosa i ziemia,
   gdy przyjdziesz sądzić świat przez ogień.}
\w{\V Tremens factus sum ego, et tímeo, dum discússio vénerit, atque ventúra ira.
   Quando caeli movéndi sunt et terra.}
  {\V Drżę cały i lękam się sądu, który nadejdzie, i przyszłego gniewu.
   Gdy poruszą się niebiosa i ziemia.}
\w{\V Dies illa, dies irae, calamitátis et misériae, dies magna et amára valde.
   Dum véneris judicáre saeculum per ignem.}
  {\V Dzień ów, dzień gniewu, klęski i nędzy, dzień wielki i bardzo gorzki.
   Gdy przyjdziesz sądzić świat przez ogień.}
\w{\V Réquiem aetérnam dona eis, Dómine: et lux perpétua lúceat eis.}
  {\V Wieczny odpoczynek racz im dać, Panie, a światłość wiekuista niechaj im świeci.}
\w{\rubr{Powtarza się} Líbera me \rubr{aż do pierwszego wersetu.}}
  {\rubr{Powtarza się} Wybaw mnie \rubr{aż do pierwszego wersetu.}}
\end{tekst}

{\centering\small\rubr{Następnie kapłan mówi: Kýrie, eléison. Christe, eléison. Kýrie,
eléison. Pater noster (po cichu), kropi i okadza trumnę, po czym odmawia modlitwę.}\par}

\clearpage
% ---------------------------------------------------------------------------
\chant{In paradisum}{antiphona · tonus VII}
{\centering\small\rubr{W drodze na cmentarz lub przy wyprowadzeniu ciała z kościoła.}\par}
\lilypondfile{../missa-pro-defunctis/scores/paradisum.ly}
\begin{tekst}
\w{In paradísum dedúcant te Angeli: in tuo advéntu suscípiant te Mártyres,
   et perdúcant te in civitátem sanctam Jerúsalem.}
  {Niech aniołowie zawiodą cię do raju, niech męczennicy wyjdą na twoje spotkanie
   i wprowadzą cię do świętego miasta Jeruzalem.}
\w{Chorus Angelórum te suscípiat, et cum Lázaro quondam páupere aetérnam hábeas réquiem.}
  {Niech cię przyjmie chór aniołów, abyś razem z Łazarzem, niegdyś ubogim,
   miał wieczny odpoczynek.}
\end{tekst}
